% Бүлэг 2

\chapter{... системийн шинжилгээ} % Бүлгийн нэр
\label{Chapter2} % Энэ бүлэг рүү ишлэл хийх бол \ref{Chapter1} командыг ашигла 

\section{Системийн үйл ажиллагааны тухай дэлгэрэнгүй}

% user scenario oruulj bolno
\section{Системийг ашиглах хэрэглэгчид}

\section{Функцийн шаардлага}

\begin{enumerate}
  \item Бүртгэл үүсгэх, нэвтрэх/гарах, нууц үг сэргээх
  \item Профайл мэдээлэл харах/засварлах
  \item Нүүр хуудсанд өөрийн бүх хичээлийг жагсаалтаар харах
  \item Хичээл бүрийн үзсэн статус, ахиц хувь болон үүсгэсэн огноо харах
  \item Шинэ хичээл үүсгэх (хоосон хэлбэрээр)
  \item Номын файл (PDF) оруулж хичээл автоматаар үүсгэх
  \item Оруулсан номыг PDF Parser-р боловсруулж хичээл→бүлэг→дэд сэдэв бүтэц гаргах
  \item Боловсруулсан бүтэц дотроос хичээлдээ оруулах сэдвийг сонгох
  \item Хичээлийн дэлгэрэнгүй хуудас: тэмдэглэлүүд, AI chatbot-ын цэс харах
  \item Хичээл дотор шинэ тэмдэглэл үүсгэх
  \item Тэмдэглэлд файл хавсаргах 
  \item Тэмдэглэлийн материал илгээмэгц AI Model API-аар сэдвийн хураангуй үүсгэх
  \item AI-аар flashcard-ууд үүсгэх
  \item AI-аар түлхүүр ойлголтууд гаргах
  \item AI-аар шалгалтын асуултууд гаргах
  \item Шалгалт өгөх
  \item Шалгалтын үр дүн 90\%-с дээш бол тэмдэглэлийг үзэж дуусгасанаар тэмдэглэх
  \item Шалгалтын дараах оноо тооцох, зөв/буруу үр дүнг харуулах
  \item Хичээлийн бүх тэмдэглэлийг үзэж дуусгасан үед хичээлийг үзэж дуусгасанаар тэмдэглэх
  \item Хэрэглэгч AI үүсгэсэн контентыг засварлах, нэмэлт чөлөөт текст бичих боломжтой байх
  \item Flashcard-уудыг нэмэх/засварлах/устгах
  \item Хичээлийн chatbot-д тухайн хичээлийн агуулгад тулгуурласан асуулт асуух
  \item Хичээлийн ахиц хувь автоматаар шинэчлэх
  \item Хичээлийг Completed гэж гараар тэмдэглэх эсвэл 100\% болсон үед автоматаар Completed болгох
  \item Хичээл, тэмдэглэлийг хайх/шүүх (нэр, огноо, ахиц, статус)
  \item Хичээлийг Public/Private болгох тохиргоо
  \item Public болгосон хичээлд хуваалцах линк үүсгэх
  \item Public хичээлийг бусад хэрэглэгч линкээр үзэх (зөвхөн унших эрх)
  \item Хэрэглэгчийн хяналтын самбар: шалгалтын амжилтын хувь, сүүлд уншсан хичээл, статусуудыг харах
  \item Админ нэвтэрч Admin Dashboard руу шилжих
  \item Админ: хэрэглэгчдийн жагсаалт, бүртгэлийн мэдээлэл, үүсгэсэн хичээлийн тоо, сүүлд нэвтэрсэн огноо харах
  \item Админ: хэрэглэгчийн хичээлүүдийн бүтэц, тэмдэглэл, AI боловсруулалтын үр дүнг харах
  \item Админ: тэмдэглэлийн материалыг дахин боловсруулуулж (reprocess) AI контентыг шинэчлэх
  \item Админ: системийн тайлан, статистик (идэвхтэй хэрэглэгч, нийт хичээл, боловсруулалтын тоо) харах
\end{enumerate}

\section{Функцийн бус шаардлага}
\begin{enumerate}
    \item Систем хэрэглэгчийн хүсэлтэд 3 секундээс багатай хугацаанд хариу өгөх.
    \item AI боловсруулалтын үр дүнд хураангуй, flashcard, тест үүсгэх үйл явц 30 секундээс хэтрэхгүй хугацаанд дууссан байх.
    \item Хэрэглэгчийн нууц үгийг шифрлэн хадгалах.
    \item Хэрэглэгчийн өгөгдөлд зөвхөн тухайн хэрэглэгч өөрөө хандах эрхтэй байх, бусад хэрэглэгчдэд харагдахгүй байх.
    \item Системийн лог бүртгэл тодорхой, алдаа илрүүлэхэд хялбар байх.
    \item Огноо/цаг: хэрэглэгчийн бүсийн \textbf{timezone}-оор зөв харуулах
    \item Систем нь HTTPS протокол ашиглан бүх харилцааг шифрлэж, мэдээллийн аюулгүй байдлыг хангах ёстой.
    \item Role-based authentication хэрэгжүүлж, хэрэглэгч, админы эрхийг ялгах ёстой.
    \item Вэб интерфэйс нь responsive дизайнтай байж, desktop болон mobile төхөөрөмж дээр адилхан ажиллагаатай байх ёстой.
    \item Систем нь алдааны болон анхааруулгын мэдэгдлийг хэрэглэгчид ойлгомжтой байдлаар харуулах ёстой.
\end{enumerate}

\section{Юзкейс диаграмм}
\begin{itemize}
    \item Actor:
    \begin{itemize}
        \item Хэрэглэгч 
        \item Админ
        \item AI Модел
    \end{itemize}
    \item Use cases:
    \begin{itemize}
        \item Нэвтрэх
        \item Хичээлүүдийг харах
        \item Хичээл үүсгэх
        \item Хичээлийн дэлгэрэнгүйг харах
        \item Тэмдэглэл үүсгэх
        \item Файл хавсаргах
        \item Тэмдэглэлийн материал боловсруулах
        \item Тэмдэглэлийн дэлгэрэнгүй агуулга харах
        \item Тэмдэглэлийн дэлгэрэнгүй агуулга засварлах
        \item Чатботоос асуулт асуух
        \item Шалгалт өгөх
        \item Шалгалтын үр дүн харах
        \item Ахиц харах
        \item Хичээлийн статусыг Completed болгох
        \item Хичээлийг нууцлал өөрчлөх
        \item Dashboard харах
        \item Хэрэглэгчдийн жагсаалт харах
        \item Хичээлүүдийн жагсаалт харах
        \item Тэмдэглэлийн AI боловсруулалтыг дахин эхлүүлэх
    \end{itemize}
\end{itemize}

% zasah
\begin{figure}[H]
    \centering
    \includegraphics[width=0.9\linewidth]{figures/chapter2/usecase.jpg}
    \caption{Use case diagram}
    \label{fig:usecase}
\end{figure}

\section{Юзкейсийн тодорхойлолт}

% iteration1 ,2  geed oruulj bolno
\section{Шинжилгээний класс диаграмм}


\section{Шинжилгээний дарааллын диаграмм}


\section{Үйл ажиллагааны диаграмм}

\section{Бүлгийн дүгнэлт}

latexmk -pdf '-interaction=nonstopmode' '-synctex=1' main.tex